
\subsection{非消失定理：定理 6.1.9 的证明}
\label{subsec:6-4-3}

对每一对
$(\boldsymbol b,b_i)
\in\mathcal M_{\mathrm{weak,def}}(L)$，
可定义上边界算子
\[
\delta^+_{(\boldsymbol b,b_1),(\boldsymbol b,b_0)}(x)
=\mathfrak q^+(e^{\boldsymbol b},e^{b_1}xe^{b_0})
\]
（见定义~\ref{def:3-8-58}
与引理~\ref{lem:3-8-60}）。
假设
$\mathfrak{PO}(b_1)=\mathfrak{PO}(b_0)$。

先证明 (6.1.10.1)。
回忆第 3.3 节的分解
\[
C(L;\Lambda_{0,\mathrm{nov}})^+
=C(L;\Lambda_{0,\mathrm{nov}})
\oplus\Lambda_{0,\mathrm{nov}}\mathbf e^+
\oplus\Lambda_{0,\mathrm{nov}}\mathbf f,
\]
其中
$\mathbf e=\operatorname{PD}([L])$
是 $(C(L;\Lambda_{0,\mathrm{nov}}),\mathfrak m)$
的同伦幺元。
\pagebreak
我们证明如下引理。

\begin{lemma}{6.4.13}\label{lem:6-4-13}
\[
\delta^+_{(\boldsymbol b,b_1),(\boldsymbol b,b_0)}(\mathbf f)
=\mathbf e^+-\mathbf e
+\mathfrak h_2^{+M}
\bigl(e^{\boldsymbol b},
e^{b_1}\mathbin{\overline\otimes}e^{b_0}\bigr).
\]
\end{lemma}

\begin{proof}
由定义~\ref{def:3-8-38}，
得到经 $\boldsymbol b$ 形变的滤过
$A_\infty$ 代数
$(C(L;\Lambda_{0,\mathrm{nov}})^+,
\mathfrak m^{+,\boldsymbol b})$。
这时 (3.8.35.4) 给出
\[
\mathfrak q_{*,1}^+(e^{\boldsymbol b},\mathbf f)
=\mathbf e^+-\mathbf e
+\mathfrak h_2^{+M}
\bigl(e^{\boldsymbol b},
1\mathbin{\overline\otimes}1\bigr).
\]
这是 (3.3.5.2) 的 ``def'' 版本。
于是由 (3.8.35.3)，
\[
\begin{aligned}
&\delta^+_{(\boldsymbol b,b_1),(\boldsymbol b,b_0)}(\mathbf f)\\
&\quad=\mathfrak q^+
\bigl(e^{\boldsymbol b},e^{b_1}\mathbf f e^{b_0}\bigr)\\
&\quad=\mathfrak q_{*,1}^+(e^{\boldsymbol b},\mathbf f)
+\mathfrak h_2^{+M}
\bigl(e^{\boldsymbol b},
e^{b_1}\mathbin{\overline\otimes}e^{b_0}\bigr)
-\mathfrak h_2^{+M}
\bigl(e^{\boldsymbol b},
1\mathbin{\overline\otimes}1\bigr)\\
&\quad=\mathbf e^+-\mathbf e
+\mathfrak h_2^{+M}
\bigl(e^{\boldsymbol b},
e^{b_1}\mathbin{\overline\otimes}e^{b_0}\bigr).
\end{aligned}
\]
\end{proof}

又因为
$\delta^+_{(\boldsymbol b,b_1),(\boldsymbol b,b_0)}
(\mathbf e^+)=0$，
便有
\[
\delta^+_{(\boldsymbol b,b_1),(\boldsymbol b,b_0)}
\Bigl(
\mathbf e-
\mathfrak h_2^{+M}
\bigl(e^{\boldsymbol b},
e^{b_1}\mathbin{\overline\otimes}e^{b_0}\bigr)
\Bigr)=0.
\]
这就证明了 (6.1.10.1)。

接下来证明 (6.1.10.2)。
由于 $\Lambda_{\mathrm{nov}}$ 系数下的结论
强于 $\Lambda_{0,\mathrm{nov}}$ 系数下的结论，
我们在前者上证明。
这里使用算子 $\mathfrak p$，\translatornote{原书实体页 387 此处从
$\operatorname{PD}([pt])\in\operatorname{Im}\delta_r$
直接写严格等式
$T^\lambda\operatorname{PD}([pt])
=\delta^+_{(\boldsymbol b,b)}(x)$。
$\delta_r$ 只作用于谱序列的滤过商页，
该前提至多给选取代表后的首项同余；
下述论证只使用此首项，
但源书没有补出代表选择的细节，见疑误 A-0376。}
并取 $b=b_1=b_0$。
反设
$\operatorname{PD}([pt])
\in\operatorname{Im}(\delta_r)$。
则可选 $\lambda\ge0$ 与
$x\in C(L;\Lambda_{0,\mathrm{nov}})$，
使相应首项满足
\[
T^\lambda\operatorname{PD}([pt])
\equiv\delta^+_{(\boldsymbol b,b)}(x)
\pmod{T^\lambda\Lambda_{0,\mathrm{nov}}^+}.
\]
由引理~\ref{lem:6-4-5}，
\[
\begin{aligned}
T^\lambda i_!\bigl(\operatorname{PD}([pt])\bigr)
&\equiv
T^\lambda\mathfrak p_1^{+\boldsymbol b,b}
\bigl(\operatorname{PD}([pt])\bigr)
\pmod{T^\lambda\Lambda_{0,\mathrm{nov}}^+}\\
&\equiv
\mathfrak p_1^{+\boldsymbol b,b}
\bigl(\delta^+_{(\boldsymbol b,b)}(x)\bigr)
=-\delta_M\bigl(
\mathfrak p_1^{+\boldsymbol b,b}(x)
\bigr)
\pmod{T^\lambda\Lambda_{0,\mathrm{nov}}^+}.
\end{aligned}
\]
因此 $i_!(\operatorname{PD}([pt]))$
在 $H^{2n}(M;\mathbb Q)$ 中为零；
这里用到 $\delta_M$ 是不含 $T$ 的通常边界算子。
但所选点类 $\operatorname{PD}([pt])\ne0$，
且 Gysin 同态把它送到 $M$ 中非零的点类，
矛盾。
\translatornote{原书实体页 387 把最后一步归因于
$i_!:H^n(L;\mathbb Q)\to H^{2n}(M;\mathbb Q)$
在整个顶次上同调上单射。
此处未假设 $L$ 连通；若 $L$ 有多个分支，
该全映射未必单射。
证明只需所选点类的 Gysin 像非零，
见疑误 A-0383。}
于是 (6.1.10.2) 得证。

现在仍取 $b_1=b_0=b$，证明 (6.1.12)。
这里使用马斯洛夫指标假设，
并用到 $H^\ell(L;\mathbb Q)$
仅在 $\ell\in\{0,\ldots,n\}$
可能非零。
注意，$C^\ell(L;\mathbb Q)
\subset S_{n-\ell}(L;\mathbb Q)$
在 $\ell<0$ 时仍可能非零：
$L$ 上可以有维数超过
$n=\dim L$ 的非零奇异链。
因此我们借助典范模型证明 (6.1.12)。

考虑
\begin{equation}
\boldsymbol b
=\sum_{i=1}^{\infty}
e^{m_i/2}T^{\lambda_i}\boldsymbol b_i
\in H^2(M;\Lambda_{0,\mathrm{nov}}^+),
\tag{6.4.14}\label{eq:6-4-14}
\end{equation}
其中 $\boldsymbol b_i\in H(M;\mathbb Q)$。
先假设 $\deg\boldsymbol b_i\ne0$；
稍后将去掉这一假设。
由总次数为 $2$ 及 $e$ 的偶次分次，有
\begin{equation}
m_i=2-\deg\boldsymbol b_i\le0.
\tag{6.4.15}\label{eq:6-4-15}
\end{equation}
置
\[
\begin{aligned}
G(L)
&=\bigl\{(\omega(\beta),\mu(\beta))
\in\mathbb R\times2\mathbb Z
\mid\mathcal M_1(\beta)\ne\varnothing\bigr\},\\
G(L;\boldsymbol b)
&=\bigl\{(\omega(\beta)+\lambda_j,
\mu(\beta)+m_j)\ \big|\\
&\hspace{18mm}{}(\omega(\beta),\mu(\beta))\in G(L),
\ j=1,2,\ldots\bigr\}
\subset\mathbb R\times2\mathbb Z,
\end{aligned}
\]
其中 $m_j,\lambda_j$ 如 (6.4.14)。

考虑滤过 $A_\infty$ 代数
$(C(L;\Lambda_{0,\mathrm{nov}})^+,
\mathfrak m^{+,\boldsymbol b})$。
它含幺，并在 $G(L;\boldsymbol b)$ 上有隙。
取其典范模型，记为
$(H(L;\Lambda_{0,\mathrm{nov}}),
\mathfrak m^{\mathrm{can},\boldsymbol b})$；\translatornote{原书实体页 387--388 的
$G(L)$、$G(L;\boldsymbol b)$
及下式 (6.4.21) 只列单个圆盘类与单个形变项，
却直接把它们用作有隙支撑集合。
按定义 3.1.6，严格的 $G$ 还须含零并对加法封闭；
$e^{\boldsymbol b}$ 与 $e^b$
的多次插入亦产生能量和。
这里保留源书集合式，
所需闭包或闭合证明未在本节给出，见疑误 A-0377。}
它也含幺，并在 $G(L;\boldsymbol b)$ 上有隙。
存在含幺且在 $G(L;\boldsymbol b)$ 上有隙的同伦等价
\[
\mathfrak f:
\bigl(H(L;\Lambda_{0,\mathrm{nov}}),
\mathfrak m^{\mathrm{can},\boldsymbol b}\bigr)
\longrightarrow
\bigl(C(L;\Lambda_{0,\mathrm{nov}})^+,
\mathfrak m^{+,\boldsymbol b}\bigr).
\]
它诱导映射
\[
\begin{aligned}
\mathfrak f_*:\quad
&\widehat{\mathcal M}_{\mathrm{weak}}
\bigl(H(L;\Lambda_{0,\mathrm{nov}}),
\mathfrak m^{\mathrm{can},\boldsymbol b}\bigr)\\
&\longrightarrow
\widehat{\mathcal M}_{\mathrm{weak}}
\bigl(C(L;\Lambda_{0,\mathrm{nov}})^+,
\mathfrak m^{+,\boldsymbol b}\bigr),
\end{aligned}
\]
并在规范等价类之间诱导同构。
对链模型中的
$b\in\widehat{\mathcal M}_{\mathrm{weak}}
(C(L;\Lambda_{0,\mathrm{nov}})^+,
\mathfrak m^{+,\boldsymbol b})$，
考虑以下条件。

\begin{statement}{条件}{6.4.16}\label{cond:6-4-16}
\begin{equation}
b=\sum_{i=1}^{\infty}
e^{k_i/2}T^{\lambda_i}b_i,
\tag{6.4.17}\label{eq:6-4-17}
\end{equation}
其中 $b_i\in C(L;\mathbb Q)^+$，
$k_i\in2\mathbb Z_{\le0}$。
\end{statement}

\begin{lemma}{6.4.18}\label{lem:6-4-18}
假设 $G(L)$ 中每个
$(\lambda_i,n_i)$ 都满足 $n_i\le0$。
则下列空间中的任意元素：
\[
\widehat{\mathcal M}_{\mathrm{weak}}
\bigl(C(L;\Lambda_{0,\mathrm{nov}})^+,
\mathfrak m^{+,\boldsymbol b}\bigr),
\]
都与某个满足
条件~\ref{cond:6-4-16} 的元素规范等价。
\end{lemma}

\begin{proof}
由假设及 (6.4.15)，
$G(L;\boldsymbol b)$ 中每个
$(\lambda,m)$ 都有 $m\le0$。
因为 $\mathfrak f$ 在该集合上有隙，
只需在典范模型的
$\widehat{\mathcal M}_{\mathrm{weak}}
(H(L;\Lambda_{0,\mathrm{nov}}),
\mathfrak m^{\mathrm{can},\boldsymbol b})$
中证明相同结论。
设
\begin{equation}
\widetilde b
=\sum_{i=1}^{\infty}
e^{k_i/2}T^{\lambda'_i}\widetilde b_i,
\tag{6.4.19}\label{eq:6-4-19}
\end{equation}
其中 $\widetilde b_i\in H(L;\mathbb Q)$。
由于典范模型中的弱定界余链属于
$H^1(L;\Lambda_{0,\mathrm{nov}})$，
有 $\deg\widetilde b_i=1-k_i$。
\translatornote{原书实体页 388
在 (6.4.19) 后把此处的
$\deg\widetilde b_i$ 写成 $\deg b_i$；
该次数式来自典范模型的展开，
见疑误 A-0378。}
而 $H^\ell(L;\mathbb Q)$
仅在 $\ell\in\{0,\ldots,n\}$
可能非零，且 $k_i$ 为偶数。
因此 $k_i\le0$，如所需。
\end{proof}

取
$\widetilde b\in
\widehat{\mathcal M}_{\mathrm{weak}}
(H(L;\Lambda_{0,\mathrm{nov}}),
\mathfrak m^{\mathrm{can},\boldsymbol b})$，
并令 $b=\mathfrak f_*(\widetilde b)$。
于是 $b$ 满足条件~\ref{cond:6-4-16}，
且任意弱定界余链都与这样的 $b$ 规范等价。
写
\begin{equation}
b=\sum_{i=1}^{\infty}
e^{n_i/2}T^{\lambda'_i}b_i,
\tag{6.4.20}\label{eq:6-4-20}
\end{equation}
并定义
\begin{equation}
\begin{aligned}
G(L;\boldsymbol b,b)
=\bigl\{&
(\omega(\beta)+\lambda'_i,
\mu(\beta)+n_i)
\in\mathbb R\times2\mathbb Z\ \big|\\
&(\omega(\beta),\mu(\beta))
\in G(L;\boldsymbol b),\quad i=1,2,\ldots
\bigr\}.
\end{aligned}
\tag{6.4.21}\label{eq:6-4-21}
\end{equation}
分别用 $\widetilde b$ 和 $b$
形变 $\mathfrak m^{\mathrm{can},\boldsymbol b}$、
$\mathfrak m^{+,\boldsymbol b}$，
得到 $\mathfrak m^{\mathrm{can},\boldsymbol b,\widetilde b}$
与 $\mathfrak m^{+,\boldsymbol b,b}$。
于是
\[
\bigl(H(L;\Lambda_{0,\mathrm{nov}}),
\mathfrak m^{\mathrm{can},\boldsymbol b,\widetilde b}\bigr),
\qquad
\bigl(C(L;\Lambda_{0,\mathrm{nov}})^+,
\mathfrak m^{+,\boldsymbol b,b}\bigr)
\]
都是在 $G(L;\boldsymbol b,b)$ 上有隙、
含幺的滤过 $A_\infty$ 代数，
且各自的 $\mathfrak m_0(1)$
均与幺元成比例。
置
\[
\delta^{\mathrm{can}}_{(\boldsymbol b,\widetilde b)}
=\mathfrak m_1^{\mathrm{can},\boldsymbol b,\widetilde b},
\qquad
\delta^+_{(\boldsymbol b,b)}
=\mathfrak m_1^{+,\boldsymbol b,b}.
\]
于是
\[
\bigl(\delta^{\mathrm{can}}_{(\boldsymbol b,\widetilde b)}\bigr)^2=0,
\qquad
\bigl(\delta^+_{(\boldsymbol b,b)}\bigr)^2=0.
\]
定义
\[
\begin{aligned}
\varphi:\quad
\bigl(H(L;\Lambda_{0,\mathrm{nov}}),
\delta^{\mathrm{can}}_{(\boldsymbol b,\widetilde b)}\bigr)
\longrightarrow
\bigl(C(L;\Lambda_{0,\mathrm{nov}})^+,
\delta^+_{(\boldsymbol b,b)}\bigr)
\end{aligned}
\]
为
\[
\varphi(x)
=\sum_{k_1,k_2}
\mathfrak f_{k_1+k_2+1}
\bigl(\widetilde b^{\otimes k_1},
x,\widetilde b^{\otimes k_2}\bigr).
\]
容易看出，$\varphi$ 是在上同调上诱导同构的链映射，\translatornote{原书实体页 388 在
$\varphi$ 的靶上误写
$\delta^{\mathrm{can}}_{(\boldsymbol b,\widetilde b)}$；
它作用于左侧典范模型，
右侧链模型微分已定义为
$\delta^+_{(\boldsymbol b,b)}$，
见疑误 A-0379。}
并在 $G(L;\boldsymbol b,b)$ 上有隙。
由 $\mathfrak f$ 的含幺性，
\begin{equation}
\varphi\bigl([\operatorname{PD}([L])]\bigr)
=\mathbf e^+.
\tag{6.4.22}\label{eq:6-4-22}
\end{equation}
这里 $[\operatorname{PD}([L])]$
是典范模型中的幺元。

现在完成定理~\ref{thm:6-1-9} 的证明。
$C(L;\Lambda_{0,\mathrm{nov}})^+$
中的 $\operatorname{PD}([pt])$
具有上同调次数 $n$，
而 $C^\ell(L;\Lambda_{0,\mathrm{nov}})^+=0$
对 $\ell>n$ 成立。
将微分展开为
\[
\delta^+_{(\boldsymbol b,b)}
=\sum_{(\lambda_i,n_i)\in G(L;\boldsymbol b,b)}
\delta^+_{(\boldsymbol b,b),\beta}
T^{\lambda_i}e^{n_i/2},
\]
其中
$\delta^+_{(\boldsymbol b,b),\beta}:
C(L;\mathbb Q)^+\to C(L;\mathbb Q)^+$
的次数为 $1-n_i$。
\translatornote{原书实体页 389
把该分量的源与靶都写成
$C(L;\mathbb Q)$；但这里的总微分
作用于扩张复形 $C(L;\mathbb Q)^+$，
见疑误 A-0380。}
由于 $n_i\le0$，
$\delta^+_{(\boldsymbol b,b),\beta}
(\operatorname{PD}([pt]))=0$。
这证明了 (6.1.12.2)。

再证明 (6.1.12.1)。
首先证明幺元
$[\operatorname{PD}([L])]
\in H^0(L;\mathbb Q)$
不在
$\delta^{\mathrm{can}}_{(\boldsymbol b,\widetilde b)}$
的像中。
写
\[
\delta^{\mathrm{can}}_{(\boldsymbol b,\widetilde b)}
=\sum_{(\lambda_i,n_i)\in G(L;\boldsymbol b,b)}
\delta^{\mathrm{can}}_{(\boldsymbol b,\widetilde b),\beta}
T^{\lambda_i}e^{n_i/2}.
\]
其中
$\delta^{\mathrm{can}}_{(\boldsymbol b,\widetilde b),\beta}:
H(L;\mathbb Q)\to H(L;\mathbb Q)$
的次数为 $1-n_i$。\translatornote{原书实体页 389
在此和式及下句的分量下标中使用链模型余链 $b$，
但左端的典范模型微分以
$\widetilde b$ 为参数；
此处统一写 $\widetilde b$，
见疑误 A-0381。}
由于 $n_i\le0$ 且
$H^\ell(L;\mathbb Q)=0$ 对 $\ell<0$，
幺元
$[\operatorname{PD}([L])]
\in H^0(L;\mathbb Q)$
不在这些分量的像中。
所以它在
$\delta^{\mathrm{can}}_{(\boldsymbol b,\widetilde b)}$
的上同调中非零。
由 (6.4.22)，
$\mathbf e^+$ 在
$H(C(L;\Lambda_{0,\mathrm{nov}})^+,
\delta^+_{(\boldsymbol b,b)})$
中也非零。
又因为
\[
\operatorname{PD}([L])
-\mathfrak h_2^{+M}
\bigl(e^{\boldsymbol b},
e^{b_1}\mathbin{\overline\otimes}e^{b_0}\bigr)
\]
与 $\mathbf e^+$ 上同调，
故它在
$\delta^+_{(\boldsymbol b,b)}$
上同调中非零。
这证明了 (6.1.12.1)。
这两个类的线性无关性
来自上边界算子保持次数滤过。

至此一直假设 $\deg\boldsymbol b_j\ne0$。
利用 (3.8.36) 与命题~\ref{prop:3-8-54}，
可以如下去掉这一假设。
上述证明中取 $C^2(M;\Lambda_{0,\mathrm{nov}})$
的一个线性子空间代表
$H^2(M;\Lambda_{0,\mathrm{nov}})$。
可以用 (3.8.36) 定义的
$C^2(M;\Lambda_{0,\mathrm{nov}})^+$
代替 $C^2(M;\Lambda_{0,\mathrm{nov}})$。
于是由 (3.8.36.3)，
可选择
\[
H^2(M;\Lambda_{0,\mathrm{nov}})
\subset C^2(M;\Lambda_{0,\mathrm{nov}})^+,
\]
使 $\operatorname{PD}([M])$
由 $[\widetilde{\mathbf e}^+]$ 代表。
现令 $\boldsymbol b_j$ 如 (6.4.14)，
并置
\[
\boldsymbol b'
=\sum_{\deg\boldsymbol b_i\ne0}
e^{m_i/2}T^{\lambda_i}\boldsymbol b_i
\in H^2(M;\Lambda_{0,\mathrm{nov}}^+).
\]
由于
$\boldsymbol b-\boldsymbol b'
=c\operatorname{PD}([M])
=c[\widetilde{\mathbf e}^+]$，
命题~\ref{prop:3-8-54} 给出
\begin{equation}
\mathfrak q(e^{\boldsymbol b},\mathbf x)
-\mathfrak q(e^{\boldsymbol b'},\mathbf x)
=\begin{cases}
0,&\mathbf x\notin B_0C(L;\Lambda_{0,\mathrm{nov}}),\\
ce\mathbf e^+,&\mathbf x=1\in B_0C(L;\Lambda_{0,\mathrm{nov}}).
\end{cases}
\tag{6.4.23}\label{eq:6-4-23}
\end{equation}
据此可证
\[
\begin{aligned}
\widehat{\mathcal M}_{\mathrm{weak}}
\bigl(C(L;\Lambda_{0,\mathrm{nov}})^+,
\mathfrak m^{+,\boldsymbol b}\bigr)
&=\widehat{\mathcal M}_{\mathrm{weak}}
\bigl(C(L;\Lambda_{0,\mathrm{nov}})^+,
\mathfrak m^{+,\boldsymbol b'}\bigr),\\
\delta^+_{(\boldsymbol b,b)}
&=\delta^+_{(\boldsymbol b',b)}
\end{aligned}
\]
对
$b\in\widehat{\mathcal M}_{\mathrm{weak}}
(C(L;\Lambda_{0,\mathrm{nov}})^+,
\mathfrak m^{+,\boldsymbol b})$
成立。
因此总可用 $\boldsymbol b'$ 代替 $\boldsymbol b$。
定理~\ref{thm:6-1-9} 至此得证。
\hfill$\square$

注意，当 $b_1\ne b_0$ 时，
(6.1.12.2$'$) 不再成立；
注~\ref{rem:6-4-12} 的例子就是反例。

\noindent\textbf{定理 L 的证明。}
假设 $L\subset M$ 相对 spin，
$H^2(L;\mathbb Q)=0$，
且 $\mu_L=0$。
由于 $\mathcal M(L;\mathbb Q)$
的障碍落在零群 $H^2(L;\mathbb Q)$ 中，
存在定界余链
$b\in\mathcal M(L;\mathbb Q)$。
由定理~\ref{thm:D} 的 \textup{(D.1)}
及 (6.1.12.1)，
元素
\[
\operatorname{PD}([L])
-\mathfrak h_2^{+M}
\bigl(1,e^b\mathbin{\overline\otimes}e^b\bigr)
\]
确定
$HF((L,b),(L,b);\Lambda_{\mathrm{nov}})$
中的非零类。
\translatornote{原书实体页 390
在本段两次把该类写成含
$\boldsymbol b,b_1,b_0$ 的一般形变式，
但此处只取无 bulk 形变的同一个 $b$；
译文代入 $\boldsymbol b=0$、$b_1=b_0=b$，
见疑误 A-0382。}
因此对任意哈密顿微分同胚 $\psi$，
有 $\psi(L)\cap L\ne\varnothing$。
若 $L$ 与 $\psi(L)$ 横截，
因 $\mu_L=0$，
还可在 $\Lambda_{\mathrm{nov}}^{(0)}$
上定义
$HF((L,b),(L,b);\Lambda_{\mathrm{nov}}^{(0)})$
（见第 5.4.6 小节）；
换言之，整数次数有定义。
上述非零类在移位前的次数为零，
故 $\psi(L)\cap L$
中存在马斯洛夫指标为 $0$ 的交点。
定理 L 得证。
\hfill$\square$
